% Supplementary S2: failure modes (moved from appendix D.2 at the
% Session 37 Stage 3 restructure; content unchanged except the internal
% cross-reference wording).

\section{Failure modes}
\label{app:failure_modes}

Four mechanistic negatives complete the estimation story. First, filter
divergence under the fixed criterion of \S\ref{sec:methods_lae} is a
transformer-forecast phenomenon: the linear filter with encoded
observations posts zero divergences where the transformer stack posts
four, and under subsampled assimilation the transformer-forecast hybrid
collapses already at every-second-frame updates while the linear filter
degrades gracefully; every-frame pressure is load-bearing for the
nonlinear filter. Second, the reconstructive baseline's estimatability
failure in the dimension grid is an observation-side failure, verified
three ways: linear probes on its true latents are healthy at every
dimension and seed, the failure is insensitive to the observation map's
ridge strength over three orders of magnitude, and the identical protocol
succeeds for the linear and predictive families in every cell. Third, the
ensemble smoother on the forecast-derived filter degrades it: at
$N = 64$ the lagged cross-covariances are sampling noise, and the
closed-form recursion on the linear stack is the working smoother.
Fourth, the oracle-covariate negative of \S\ref{app:rex_ledger}
belongs to this list from the deployment side: more information, worse
forecast, when the information indexes a family of held-out dynamics
rather than the current state.
